\section{Proofs}\label{app:proofs}

\noindent\textbf{Lemma 1.} If $\{S_\lambda\}$ is nested and
$L_i(\lambda)=f(S_\lambda\cap Y_i)$ with $f$ non-increasing under set inclusion, then
$\lambda\mapsto L_i(\lambda)$ is non-increasing.\\
\emph{Proof.} $\lambda\le\lambda'\Rightarrow S_\lambda\cap Y_i\subseteq S_{\lambda'}\cap Y_i$, so
$f(S_{\lambda'}\cap Y_i)\le f(S_\lambda\cap Y_i)$. Each indicator in $L^{\mathrm{les}}$ and
$L^{\mathrm{small}}$ can only flip from 1 to 0 as the covered fraction of a lesion grows. $\square$

\medskip\noindent\textbf{Lemma 2.} Let $q_k$ be the $\lceil\tau|Y^k|\rceil$-th largest score in $Y^k$.
Lesion $k$ is hit at cutoff $c$ iff $q_k\ge c$.\\
\emph{Proof.} At least $\tau|Y^k|$ pixels of $Y^k$ score $\ge c$ iff the
$\lceil\tau|Y^k|\rceil$-th largest does. $\square$

\medskip\noindent\textbf{Proposition 1 (validity with a tuned family).} Let $\mathcal T$ be a finite
set of families, and let $D_{\mathrm{tune}}$ and $D_{\mathrm{cal}}$ be disjoint, with
$D_{\mathrm{cal}}\cup\{n{+}1\}$ exchangeable conditionally on $D_{\mathrm{tune}}$. For any selection
rule $\hat T=\Psi(D_{\mathrm{tune}})$ and every risk $r$,
$\mathbb E[L^r_{n+1}(\hat\lambda^\cup)]\le\alpha$.\\
\emph{Proof.} Conditionally on $D_{\mathrm{tune}}$, $\hat T$ is a fixed label-free transform, so the
sets are nested and Lemma~1 applies. Split CRC~\cite{angelopoulos2024crc} gives
$\mathbb E[L^r_{n+1}(\hat\lambda^r)\mid D_{\mathrm{tune}}]\le\alpha$ on the subpopulation where $L^r$ is
defined. Monotonicity gives $L^r_{n+1}(\hat\lambda^\cup)\le L^r_{n+1}(\hat\lambda^r)$ pointwise. Taking
the expectation over $D_{\mathrm{tune}}$ completes the proof. $\square$

\medskip\noindent\textbf{Lemma 3.} For $\rho\in[0,1]$, $s_{\rho,t_0}\in[0,1]$. Every candidate pixel has
$s\ge t_0$, every non-candidate pixel has $s<t_0$, and $s_{0,t_0}=\hat p$.\\
\emph{Proof.} On $C_j$, $t_0\le\hat p\le P_j\le 1$ and $P_j\le P_j^\rho\le 1$. Hence
$t_0\le\hat p\le\hat p/P_j^\rho\le\hat p/P_j\le 1$. Off the candidates, $s=\hat p<t_0$. $\square$

\medskip\noindent\textbf{Proposition 2 (amplitude invariance).} If $\hat p=a_jg_j$ on $C_j$ with $a_j>0$,
then $s_{1,t_0}=g_j/\max_{C_j}g_j$ on $C_j$. Hence the hit condition of every lesion inside $C_j$
(Lemma~2) does not depend on $a_j$.\\
\emph{Proof.} $P_j=a_j\max_{C_j}g_j$; substitute into~\eqref{eq:le}. $\square$

\medskip\noindent\textbf{Proposition 3 (background exclusion).} For every $\rho$ and every cutoff
$c\ge t_0$, $\{s_{\rho,t_0}\ge c\}\subseteq\bigcup_jC_j$.\\
\emph{Proof.} Immediate from Lemma~3. $\square$

\noindent Thresholding gives no such containment. If the joint constraint needs a
thresholding cutoff $c<t_0$, every pixel in $\{c\le\hat p<t_0\}$ joins the set.
The equalized family can meet the same constraint at a cutoff $\ge t_0$ whenever
every lesion that must be hit lies inside a candidate. This is a sufficient
mechanism, not an optimality result. The loss-optimal nested family for a
lesion-level risk depends on the joint law of $(\hat p,Y)$ and remains open.
